\documentclass[12pt,a4paper]{ctexbook}
\usepackage{geometry}
\geometry{margin=2.2cm}
\usepackage{setspace}
\setstretch{1.35}
\usepackage{hyperref}
\title{全宋词选·ci.song.1000}
\author{古典文献整理}
\date{}
\begin{document}
\maketitle
\tableofcontents
\newpage
\section{踏莎行}
\textbf{韩绛}\par\vspace{0.5em}
嵩峤云高，洛川波暖。\par
举头乔木森无断。\par
□□□雨绝风尘，小桥频过春渠满。\par
□□离宫，□棱斗焕。\par
万家罗绮多游伴。\par
□□□□自风□，□□是处喧弦管。\par

\section{菩萨蛮}
\textbf{李师中}\par\vspace{0.5em}
子规啼破城楼月。\par
画船晓载笙歌发。\par
两岸荔枝红。\par
万家烟雨中。\par
佳人相对泣。\par
泪下罗衣湿。\par
从此信音稀。\par
岭南无雁飞。\par

\section{喜迁莺}
\textbf{蔡挺}\par\vspace{0.5em}
霜天清晓。\par
望紫塞古垒，寒云衰草。\par
汗马嘶风，边鸿翻月，垅上铁衣寒早。\par
剑歌骑曲悲壮，尽道君恩难报。\par
塞垣乐，尽双鞬锦带，山西年少。\par
谈笑。\par
刁斗静。\par
烽火一把，常送平安耗。\par
圣主忧边，威灵遐布，骄虏且宽天讨。\par
岁华向晚愁思，谁念玉关人老。\par
太平也，且欢娱，不惜金尊频倒。\par

\section{喜长新}
\textbf{王益柔}\par\vspace{0.5em}
秋云朔吹晓徘徊。\par
雪照楼台。\par
梁王宴召有邹枚。\par
相如独逞雄才。\par
明烛薰炉香暖，深劝金杯。\par
庭前粉艳有寒梅。\par
一枝昨夜先开。\par

\section{西江月}
\textbf{韩维}\par\vspace{0.5em}
早岁相期林下，高年同在尊前。\par
风花绣舞乍晴天。\par
绿蚁新浮酒面。\par
身外虚名电转，人间急景梭传。\par
当筵莫惜听朱弦。\par
一品归来强健。\par

\section{踏莎行}
\textbf{韩维}\par\vspace{0.5em}
归雁低空，游蜂趁暖。\par
凭高目向西云断。\par
具茨山外夕阳多，展江亭下春波满。\par
双桂情深，千花明焕。\par
良辰谁是同游伴。\par
辛夷花谢早梅开。\par
应须次第调弦管。\par

\section{减字木兰花}
\textbf{韩维}\par\vspace{0.5em}
和风动□。\par
□□新年□入手。\par
世事尘□。\par
□□□情近酒□。\par
水开湖□。\par
□□笙歌波面起。\par
相与排□。\par
□□□胜特地□。\par

\section{浪淘沙}
\textbf{韩维}\par\vspace{0.5em}
饱食日□□。\par
□上危亭。\par
东风昨夜入□□。\par
□□雪晴云□□，□遍银屏。\par
回首叹劳生。\par
□鼎相承。\par
皇恩早晚□□□。\par
□□陂边垂钓手，不负幽情。\par

\section{胡捣练令・胡捣练}
\textbf{韩维}\par\vspace{0.5em}
夜来风横雨飞狂，满地闲花衰草。\par
燕子渐归春悄。\par
帘幕垂清晓。\par
天将佳景与闲人，美酒宁嫌华皓。\par
留取旧时欢笑。\par
莫共秋光老。\par

\section{失调名}
\textbf{韩维}\par\vspace{0.5em}
轻云薄雾。\par
散作催花雨。\par

\section{失调名}
\textbf{韩维}\par\vspace{0.5em}
兄弟对举杯。\par

\section{赏南枝}
\textbf{曾巩}\par\vspace{0.5em}
暮冬天地闭，正柔木冻折，瑞雪飘飞。\par
对景见南山，岭梅露、几点清雅容姿。\par
丹染萼、玉缀枝。\par
又岂是、一阳有私。\par
大抵是、化工独许，使占却先时。\par
霜威莫苦凌持。\par
此花根性，想群卉争知。\par
贵用在和羹，三春里、不管绿是红非。\par
攀赏处、宜酒卮。\par
醉拈嗅、幽香更奇。\par
倚阑干、仗何人去，嘱羌管休吹。\par

\section{阮郎归}
\textbf{司马光}\par\vspace{0.5em}
渔舟容易入春山。\par
仙家日月闲。\par
绮窗纱幌映朱颜。\par
相逢醉梦间。\par
松露冷，海霞殷。\par
匆匆整棹还。\par
落花寂寂水潺潺。\par
重寻此路难。\par

\section{西江月}
\textbf{司马光}\par\vspace{0.5em}
宝髻松松挽就，铅华淡淡妆成。\par
青烟翠雾罩轻盈。\par
飞絮游丝无定。\par
相见争如不见，有情何似无情。\par
笙歌散后酒初醒。\par
深院月斜人静。\par

\section{锦堂春・锦堂春慢}
\textbf{司马光}\par\vspace{0.5em}
红日迟迟，虚廊转影，槐阴迤逦西斜。\par
彩笔工夫，难状晚景烟霞。\par
蝶尚不知春去，漫绕幽砌寻花。\par
奈猛风过后，纵有残红，飞向谁家。\par
始知青鬓无价，叹飘零官路，荏苒年华。\par
今日笙歌丛里，特地咨嗟。\par
席上青衫湿透，算感旧、何止琵琶。\par
怎不教人易老，多少离愁，散在天涯。\par

\section{临江仙}
\textbf{苏氏}\par\vspace{0.5em}
一夜东风穿绣户，融融暖应佳时。\par
春来何处最先知。\par
平明堤上柳。\par
染遍郁金枝。\par
姊妹嬉游时节近，今朝应怨来迟。\par
凭谁说与到家期。\par
玉钗头上胜，留待远人归。\par

\section{更漏子}
\textbf{苏氏}\par\vspace{0.5em}
小阑干，深院宇。\par
依旧当时别处。\par
朱户锁，玉楼空。\par
一帘霜日红。\par
弄珠江，何处是，望断碧云无际。\par
凝泪眼，出重城。\par
隔溪羌笛声。\par

\section{鹊桥仙}
\textbf{苏氏}\par\vspace{0.5em}
星移斗转，玉蟾西下，渐觉东郊向晓。\par
马嘶人语隔霜林，望千里、长安古道。\par
珠宫姊妹，相逢方信，别后十分瘦了。\par
上林归去正花时，争奈向、花前又老。\par

\section{踏莎行}
\textbf{苏氏}\par\vspace{0.5em}
孤馆深沈，晓寒天气。\par
解鞍独自阑干倚。\par
暗香浮动月黄昏，落梅风送沾衣袂。\par
待写红笺，凭谁与寄。\par
先教觅取嬉游地。\par
到家正是早春时，小桃花下拼沈醉。\par

\section{清平乐}
\textbf{刘敞}\par\vspace{0.5em}
小山丛桂。\par
最有留人意。\par
拂叶攀花无限思。\par
雨湿浓香满袂。\par
别来过了秋光。\par
翠帘昨夜新霜。\par
多少月宫闲地，姮娥与借微芳。\par

\section{踏莎行}
\textbf{刘敞}\par\vspace{0.5em}
蜡炬高高，龙烟细细。\par
玉楼十二门初闭。\par
疏帘不卷水晶寒，小屏半掩琉璃翠。\par
桃叶新声，榴花美味。\par
南山宾客东山妓。\par
利名不肯放人闲，忙中偷取工夫醉。\par

\section{导引}
\textbf{王}\par\vspace{0.5em}
忆玉清景，繁盛极当时。\par
千古事难追。\par
汉家别庙秋风起，空出奉宸衣。\par
三山浮海日晖晖。\par
羽盖共云飞。\par
灵宫旧是栖真处，还望玉舆归。\par

\section{发引}
\textbf{王}\par\vspace{0.5em}
玉宸朝晚，忽掩赭黄衣。\par
愁雾锁金扉。\par
蓬莱待得仙丹至，人世已成非。\par
龙轩天仗转西畿。\par
旌旆入云飞。\par
望陵宫女垂红泪，不见翠舆归。\par

\section{发引}
\textbf{王}\par\vspace{0.5em}
上林春晚，曾是奉宸游。\par
水殿戏龙舟。\par
玉箫吹断催仙驭，一去隔千秋。\par
游人重到曲江头。\par
事往涕难收。\par
空馀御幄传觞处，依旧水东流。\par

\section{凤箫吟・芳草}
\textbf{韩缜}\par\vspace{0.5em}
锁离愁，连绵无际，来时陌上初熏。\par
绣帏人念远，暗垂珠泪，泣送征轮。\par
长亭长在眼，更重重、远水孤云。\par
但望极楼高，尽日目断王孙。\par
消魂。\par
池塘别后，曾行处、绿妒轻裙。\par
恁时携素手，乱花飞絮里，缓步香裀。\par
朱颜空自改，向年年、芳意长新。\par
遍绿野，嬉游醉眠，莫负青春。\par

\section{蝶恋花}
\textbf{韩缜姬}\par\vspace{0.5em}
香作风光浓著露。\par
正恁双栖，又遣分飞去。\par
密诉东君应不许。\par
泪波一洒奴衷素。\par

\section{浪淘沙}
\textbf{裴湘}\par\vspace{0.5em}
雁塞说并门。\par
郡枕西汾。\par
山形高下远相吞。\par
古寺楼台依碧嶂，烟景遥分。\par
晋庙锁溪云。\par
箫鼓仍存。\par
牛羊斜日自归村。\par
惟有故城禾黍地，前事消魂。\par

\section{浪淘沙}
\textbf{裴湘}\par\vspace{0.5em}
万国仰神京。\par
礼乐纵横。\par
葱葱佳气锁龙城。\par
日御明堂天子圣，朝会簪缨。\par
九陌六街平。\par
万国充盈。\par
青楼弦管酒如渑。\par
别有隋堤烟柳暮，千古含情。\par

\section{花心动}
\textbf{阮逸女}\par\vspace{0.5em}
仙苑春浓，小桃开，枝枝已堪攀折。\par
乍雨乍晴，轻暖轻寒，渐近赏花时节。\par
柳摇台榭东风软，帘栊静、幽禽调舌。\par
断魂远、闲寻翠径，顿成愁结。\par
此恨无人共说。\par
还立尽黄昏，寸心空切。\par
强整绣衾，独掩朱扉，簟枕为谁铺设。\par
夜长更漏传声远，纱窗映、银缸明灭。\par
梦回处，梅梢半笼淡月。\par

\section{蝶恋花}
\textbf{滕甫}\par\vspace{0.5em}
叶底无风池面静。\par
掬水佳人，拍破青铜镜。\par
残月朦胧花弄影。\par
新梳斜插乌云鬓。\par
拍索闷怀添酒兴。\par
旋撷园蔬，随分成盘饤。\par
说与翠微休急性。\par
功名富贵皆前定。\par

\section{蝶恋花}
\textbf{滕甫}\par\vspace{0.5em}
昼永无人深院静。\par
一枕春醒，犹未忺临镜。\par
帘卷新蟾光射影。\par
速忙掠起蓬松鬓。\par
对景沈吟嗟没兴。\par
薄幸不来，空把杯盘饤。\par
休道妇人多水性。\par
今宵独自言无定。\par

\section{桂枝香}
\textbf{王安石}\par\vspace{0.5em}
登临送目。\par
正故国晚秋，天气初肃。\par
千里澄江似练，翠峰如簇。\par
归帆去棹残阳里，背西风、酒旗斜矗。\par
彩舟云淡，星河鹭起，画图难足。\par
念往昔、繁华竞逐。\par
叹门外楼头，悲恨相续。\par
千古凭高，对此漫嗟荣辱。\par
六朝旧事随流水，但寒烟、芳草凝绿。\par
至今商女，时时犹唱，后庭遗曲。\par

\section{甘露歌}
\textbf{王安石}\par\vspace{0.5em}
折得一枝香在手。\par
人间应未有。\par
疑是经春雪未消。\par
今日是何朝。\par

\section{甘露歌}
\textbf{王安石}\par\vspace{0.5em}
尽日含毫难比兴。\par
都无色可并。\par
万里晴天何处来。\par
真是屑琼瑰。\par

\section{甘露歌}
\textbf{王安石}\par\vspace{0.5em}
天寒日暮山谷里。\par
的砾愁成水。\par
池上渐多枝上稀。\par
唯有故人知。\par

\section{菩萨蛮}
\textbf{王安石}\par\vspace{0.5em}
数家茅屋闲临水。\par
单衫短帽垂杨里。\par
今日是何朝。\par
看予度石桥。\par
梢梢新月偃。\par
午醉醒来晚。\par
何物最关情。\par
黄鹂三两声。\par

\section{渔家傲}
\textbf{王安石}\par\vspace{0.5em}
灯火已收正月半。\par
山南山北花撩乱。\par
闻说洊亭新水漫。\par
骑款段。\par
穿云入坞寻游伴。\par
却拂僧床褰素幔。\par
千岩万壑春风暖。\par
一弄松声悲急管。\par
吹梦断。\par
西看窗日犹嫌短。\par

\section{渔家傲}
\textbf{王安石}\par\vspace{0.5em}
平岸小桥千嶂抱。\par
柔蓝一水萦花草。\par
茅屋数间窗窈窕。\par
尘不到。\par
时时自有春风扫。\par
午枕觉来闻语鸟。\par
敧眠似听朝鸡早。\par
忽忆故人今总老。\par
贪梦好。\par
茫然忘了邯郸道。\par

\section{雨霖铃}
\textbf{王安石}\par\vspace{0.5em}
孜孜矻矻。\par
向无明里、强作窠窟。\par
浮名浮利何济，堪留恋处，轮回仓猝。\par
幸有明空妙觉，可弹指超出。\par
缘底事、抛了全潮，认一浮沤作瀛渤。\par
本源自性天真佛。\par
祗些些、妄想中埋没。\par
贪他眼花阳艳，谁信道、本来无物。\par
一旦茫然，终被阎罗老子相屈。\par
便纵有、千种机筹，怎免伊唐突。\par

\section{清平乐}
\textbf{王安石}\par\vspace{0.5em}
云垂平野。\par
掩映竹篱茅舍。\par
阒寂幽居实潇洒。\par
是处绿娇红冶。\par
丈夫运用堂堂。\par
且莫五角六张。\par
若有一卮芳酒，逍遥自在无妨。\par

\section{浣溪沙}
\textbf{王安石}\par\vspace{0.5em}
百亩中庭半是苔。\par
门前白道水萦回。\par
爱闲能有几人来。\par
小院回廊春寂寂，山桃溪杏两三栽。\par
为谁零落为谁开。\par

\section{诉衷情}
\textbf{王安石}\par\vspace{0.5em}
常时黄色见眉间。\par
松桂我同攀。\par
每言天上辛苦，不肯饵金丹。\par
怜水静，爱云闲。\par
便忘还。\par
高歌一曲，岩谷迤逦，宛似商山。\par

\section{诉衷情}
\textbf{王安石}\par\vspace{0.5em}
练巾藜杖白云间。\par
有兴即跻攀。\par
追思往昔如梦，华毂也曾丹。\par
尘自扰，性长闲。\par
更无还。\par
达如周召，穷似丘轲，祗个山山。\par

\section{诉衷情}
\textbf{王安石}\par\vspace{0.5em}
茫然不肯住林间。\par
有处即追攀。\par
将他死语图度，怎得离真丹。\par
浆水价，匹如闲。\par
也须还。\par
何如直截，踢倒军持，赢取沩山。\par

\section{诉衷情}
\textbf{王安石}\par\vspace{0.5em}
营巢燕子逞翱翔。\par
微志在雕梁。\par
碧云举翮千里，其奈有鸾皇。\par
临济处，德山行。\par
果承当。\par
自时降住，一切天魔，扫地焚香。\par

\section{诉衷情}
\textbf{王安石}\par\vspace{0.5em}
莫言普化祗颠狂。\par
真解作津梁。\par
蓦然打个筋斗，直跳过羲皇。\par
临济处，德山行。\par
果承当。\par
将他建立，认作心诚，也是寻香。\par

\section{南乡子}
\textbf{王安石}\par\vspace{0.5em}
嗟见世间人。\par
但有纤毫即是尘。\par
不住旧时无相貌，沈沦。\par
祗为从来认识神。\par
作麽有疏亲。\par
我自降魔转法轮。\par
不是摄心除妄想，求真。\par
幻化空身即法身。\par

\section{南乡子}
\textbf{王安石}\par\vspace{0.5em}
自古帝王州。\par
郁郁葱葱佳气浮。\par
四百年来成一梦，堪愁。\par
晋代衣冠成古丘。\par
绕水恣行游。\par
上尽层城更上楼。\par
往事悠悠君莫问，回头。\par
槛外长江空自流。\par

\section{浪淘沙令}
\textbf{王安石}\par\vspace{0.5em}
伊吕两衰翁。\par
历遍穷通。\par
一为钓叟一耕佣。\par
若使当时身不遇，老了英雄。\par
汤武偶相逢。\par
风虎云龙。\par
兴王祗在笑谈中。\par
直至如今千载后，谁与争功。\par

\section{望江南・忆江南}
\textbf{王安石}\par\vspace{0.5em}
归依众，梵行四威仪。\par
愿我遍游诸佛土。\par
十方贤圣不相离。\par
永灭世间痴。\par

\section{望江南・忆江南}
\textbf{王安石}\par\vspace{0.5em}
归依法，法法不思议。\par
愿我六根常寂静，心如宝月映琉璃。\par
了法更无疑。\par

\section{望江南・忆江南}
\textbf{王安石}\par\vspace{0.5em}
归依佛，弹指越三祗。\par
愿我速登无上觉，还如佛坐道场时。\par
能智又能悲。\par

\section{望江南・忆江南}
\textbf{王安石}\par\vspace{0.5em}
三界里，有取总灾危。\par
普愿众生同我愿，能於空有善思惟。\par
三宝共住持。\par

\section{西江月}
\textbf{王安石}\par\vspace{0.5em}
梅好惟嫌淡伫，天教薄与胭脂。\par
真妃初出华清池。\par
酒入琼姬半醉。\par
东阁诗情易动，高楼玉管休吹。\par
北人浑作杏花疑。\par
惟有青枝不似。\par

\section{渔家傲}
\textbf{王安石}\par\vspace{0.5em}
隔岸桃花红未半。\par
枝头已有蜂儿乱。\par
惆怅武陵人不管。\par
清梦断。\par
亭亭伫立春宵短。\par

\section{清平乐}
\textbf{王安石}\par\vspace{0.5em}
留春不住。\par
费尽莺儿语。\par
满地残红宫锦污。\par
昨夜南园风雨。\par
小怜初上琵琶。\par
晓来思绕天涯。\par
不肯画堂朱户，春风自在杨花。\par

\section{生查子}
\textbf{王安石}\par\vspace{0.5em}
雨打江南树。\par
一夜花开无数。\par
绿叶渐成阴，下有游人归路。\par
与君相逢处。\par
不道春将暮。\par
把酒祝东风，且莫恁、匆匆去。\par

\section{谒金门}
\textbf{王安石}\par\vspace{0.5em}
春又老。\par
南陌酒香梅小。\par
遍地落花浑不扫。\par
梦回情意悄。\par
红笺寄与添烦恼。\par
细写相思多少。\par
醉后几行书字小。\par
泪痕都搵了。\par

\section{菩萨蛮}
\textbf{王安石}\par\vspace{0.5em}
海棠乱发皆临水。\par
君知此处花何似。\par
凉月白纷纷。\par
香风隔岸闻。\par
啭枝黄鸟近。\par
隔岸声相应。\par
随意坐莓苔。\par
飘零酒一杯。\par

\section{千秋岁引}
\textbf{王安石}\par\vspace{0.5em}
别馆寒砧，孤城画角。\par
一派秋声入寥廓。\par
东归燕从海上去，南来雁向沙头落。\par
楚台风，庾楼月，宛如昨。\par
无奈被些名利缚。\par
无奈被他情担阁。\par
可惜风流总闲却。\par
当初漫留华表语，而今误我秦楼约。\par
梦阑时，酒醒后，思量著。\par

\section{定风波}
\textbf{吴氏1}\par\vspace{0.5em}
待得明年重把酒。\par
携手。\par
那知无雨又无风。\par

\section{蜡梅香}
\textbf{吴师孟}\par\vspace{0.5em}
锦里阳和，看万木凋时，早梅独秀。\par
珍馆琼楼畔，正降跗初吐，秾华将茂。\par
国艳天葩，真澹伫、雪肌清瘦。\par
似广寒宫，铅华未御，自然妆就。\par
凝睇倚朱阑，喷清香暗度，易袭襟袖。\par
好与花为主，宜秉烛、频观泛湘酎。\par
莫待南枝，随乐府、新声吹后。\par
对赏心人，良辰好景，须信难偶。\par

\section{阮郎归}
\textbf{俞紫芝}\par\vspace{0.5em}
钓鱼船上谢三郎。\par
双鬓已苍苍。\par
蓑衣未必清贵，不肯换金章。\par
汀草畔，浦花旁。\par
静鸣榔。\par
自来好个，渔父家风，一片潇湘。\par

\section{临江仙}
\textbf{俞紫芝}\par\vspace{0.5em}
弄水亭前千万景，登临不忍空回。\par
水轻墨澹写蓬莱。\par
莫教世眼，容易洗尘埃。\par
收去雨昏都不见，展时还似云开。\par
先生高趣更多才。\par
人人尽道，小杜却重来。\par

\section{苏幕遮}
\textbf{俞紫芝}\par\vspace{0.5em}
地锺灵，天应瑞。\par
簇簇香苞、团作真珠蕊。\par
玉宇瑶台分十二。\par
要伴姮娥，月里双双睡。\par
月如花，花似月。\par
花月生香，添此真奇异。\par
不许扬州夸间气。\par
昨夜春风，唤醒琼琼醉。\par

\section{失调名}
\textbf{郑獬}\par\vspace{0.5em}
玉环妾意无渝。\par
问君心、朝槿何如。\par

\section{好事近}
\textbf{郑獬}\par\vspace{0.5em}
江上探春回，正值早梅时节。\par
两行小槽双凤，按凉州初彻。\par
谢娘扶下绣鞍来，红靴踏残雪。\par
归去不须银烛，有山头明月。\par

\section{好事近}
\textbf{郑獬}\par\vspace{0.5em}
把酒对江梅，花小未禁风力。\par
何计不教零落，为青春留得。\par
故人莫问在天涯，尊前苦相忆。\par
好把素香收取，寄江南消息。\par

\section{渔家傲}
\textbf{强至}\par\vspace{0.5em}
雪月照梅溪畔路。\par
幽姿背立无言语。\par
冷浸瘦枝清浅处。\par
香暗度。\par
妆成处士横斜句。\par
浑似玉人常淡伫。\par
菱花相对盈清楚。\par
谁解小图先画取。\par
天欲曙。\par
恐随月色云间去。\par

\section{踏莎行}
\textbf{沈注}\par\vspace{0.5em}
竹阁云深，巢虚人阒。\par
几年湖上音尘寂。\par
风流今有使君家，月明夜夜闻双笛。\par

\section{望梅花}
\textbf{蒲宗孟}\par\vspace{0.5em}
一阳初起。\par
暖力未胜寒气。\par
堪赏素华长独秀，不并开红抽紫。\par
青帝只应怜洁白，不使雷同众卉。\par
淡然难比。\par
粉蝶岂知芳蕊。\par
半夜卷帘如乍失，只在银蟾影里。\par
残雪枝头君认取，自有清香旖旎。\par

\section{减字木兰花}
\textbf{陈汝羲}\par\vspace{0.5em}
纤纤素手。\par
盘里醉花新点就。\par
对叶双心。\par
别有东风意思深。\par
琼沾粉缀。\par
消得玉堂留客醉。\par
试嗅清芳。\par
别有红罗巧袖香。\par

\section{行香子}
\textbf{汪辅之}\par\vspace{0.5em}
晚绿寒红。\par
芳意匆匆。\par
惜年华、今与谁同。\par
碧云零落，数字宾鸿。\par
看渚莲凋，宫扇旧，怨秋风。\par
流波坠叶，佳期何在，想天教、离恨无穷。\par
试将前事，闲倚梧桐。\par
有销魂处，明月夜，锦屏空。\par

\section{鹧鸪天}
\textbf{范纯仁}\par\vspace{0.5em}
腊后春前暖律催。\par
日和风软欲开梅。\par
公方结客寻佳景，我亦忘形趁酒杯。\par
添歌管，续尊罍。\par
更阑烛短未能回。\par
清欢莫待相期约，乘兴来时便可来。\par

\section{雨中花・夜行船}
\textbf{张才翁}\par\vspace{0.5em}
万缕青青，初眠官柳，向人犹未成阴。\par
据雕鞍马上，拥鼻微吟。\par
远宦情怀谁问，空嗟壮志销沈。\par
正好花时节，山城留滞、忍负归心。\par
别离万里，飘蓬无定，谁念会合难凭。\par
相聚里，休辞金盏，酒浅还深。\par
欲把春愁抖擞，春愁转更难禁。\par
乱山高处，凭阑垂袖，聊寄登临。\par

\section{渔家傲}
\textbf{寿涯禅师}\par\vspace{0.5em}
深愿弘慈无缝罅。\par
乘时走入众生界。\par
窈窕丰姿都没赛。\par
提鱼卖。\par
堪笑马郎来纳败。\par
清冷露湿金襽坏。\par
茜裙不把珠缨盖。\par
特地掀来呈捏怪。\par
牵人爱。\par
还尽许多菩萨债。\par

\section{水龙吟}
\textbf{章}\par\vspace{0.5em}
燕忙莺懒花残，正堤上、柳花飘坠。\par
轻飞点画青林，谁道全无才思。\par
闲趁游丝，静临深院，日长门闭。\par
傍珠帘散漫，垂垂欲下，依前被，风扶起。\par
兰帐玉人睡觉，怪春衣、雪粘琼缀。\par
绣床旋满，香球无数，才圆却碎。\par
时见蜂儿，仰粘轻粉，鱼吹池水。\par
望章台路杳，金鞍游荡，有盈盈泪。\par

\section{声声令・胜胜令}
\textbf{章}\par\vspace{0.5em}
帘移碎影，香褪衣襟。\par
旧家庭院嫩苔侵。\par
东风过尽，暮云锁，绿窗深。\par
怕对人、闲枕剩衾。\par
楼底轻阴。\par
春信断，怯登临。\par
断肠魂梦两沈沈。\par
花飞水远，便从今。\par
莫追寻。\par
又怎禁、蓦地上心。\par

\section{渔父乐・渔歌子}
\textbf{徐积}\par\vspace{0.5em}
水曲山隈四五家。\par
夕阳烟火隔芦花。\par
渔唱歇，醉眠斜。\par
纶竿蓑笠是生涯。\par

\section{渔父乐・渔歌子}
\textbf{徐积}\par\vspace{0.5em}
见说红尘罩九衢。\par
贪名逐利各区区。\par
论得失，问荣枯。\par
争似侬家占五湖。\par

\end{document}